\subsection{结果段落模板：按 claim--evidence--statistic 组织首批结果}\label{subsec:typed-address-biaxial-completion-window6-results-paragraph-templates}
本节不再引入新对象，而是把上一节所固定的首批图表与表格压成可直接并入正文的结果段落模板。
其写法遵循一个简单而严格的原则：每段只服务于一个主张，先给对象事实，再给图表证据，最后给数值证书；至于机制解释与物理比喻，则尽量后移到讨论层。
如此安排后，读者会先看到 window-\(6\) 的刚性离散事实，再看到这些事实如何支撑粗糙可见量、表观随机、boundary holonomy 与账本预算四条主线的首次汇合。

\subsubsection{组织原则}\label{subsubsec:typed-address-biaxial-completion-window6-results-organization-principle}
\begin{proposition}[首批结果段的组织原则]\label{prop:typed-address-biaxial-completion-window6-results-organization-principle}
\leanverified{paper\_typed\_address\_biaxial\_completion\_window6\_results\_organization\_principle}
围绕 window-\(6\) 的首批结果段，应统一采用如下结构：
\begin{enumerate}
  \item 以一句 claim 固定该段唯一主张；
  \item 以一句 evidence 指向上一节已经锁定的对象图、表或闭式证书；
  \item 以一句 statistic 给出最小而充分的精确常数、直方图或位数阈值；
  \item 仅在前三句完成后，才允许补充机制解释。
\end{enumerate}
特别地，若某一段所依赖的图表仍未通过上一节的协议门，则该段只能以占位模板出现，而不能被写成既成结果。
\end{proposition}

\begin{proof}
这只是把 \S\ref{subsec:typed-address-biaxial-completion-window6-numerical-audit-protocol} 与 \S\ref{subsec:typed-address-biaxial-completion-window6-figure-table-guidelines} 的顺序约束改写成正文写作规则。
在 fixed-window 闭式、boundary holonomy 与预算裂分未锁定之前，任何统计或几何解释都缺少对象层宿主，因此只能作为未获准的外推而不能进入结果段。证毕。
\end{proof}

\subsubsection{可见层锚点：window-\texorpdfstring{$6$}{6} 已经是首个非平凡多壳实例}\label{subsubsec:typed-address-biaxial-completion-window6-results-visible-anchor}
围绕可见层锚点的首段结果，可写成如下形式。

window-\(6\) 已经形成首个非平凡的可见锚点。
具体地，bin-fold 将 \(64\) 个微态压缩到 \(21\) 个稳定类型，其中 cyclic 扇区与 boundary 扇区分别含 \(18\) 与 \(3\) 个词；对应纤维大小恰落在 \(\{2,3,4\}\) 三层，刚性直方图为 \(2{:}8,3{:}4,4{:}9\)。
配合上一节的可见层与纤维谱锚点图及刚性不变量总表，可见层的多壳结构与规范体积锚点同步锁定；其中
\[
\log|G_6^{\mathrm{bin}}|=39\log 2+13\log 3\approx 41.314700,
\qquad
g_6\approx 0.645542,
\qquad
\kappa_6\approx 1.159067,
\]
并且
\[
\dim_{\FF_2}\bigl(G_6^{\mathrm{bin}}\bigr)^{\mathrm{ab}}=21.
\]
因此，最小窗口已经不是单壳 coarse-graining，而是一个同时携带多壳纤维残差、规范体积与异常荷格的首个可审计实例。

这一段的功能，是把后文一切关于粗糙可见性、表观随机、boundary holonomy 与账本预算的讨论都钉在一个固定的有限窗硬对象上，而不是漂浮在解释层。

\subsubsection{boundary holonomy：最小曲率候选已经在三条 boundary 纤维上出现}\label{subsubsec:typed-address-biaxial-completion-window6-results-boundary-holonomy}
围绕 boundary holonomy 的结果段，可写成如下形式。

boundary 扇区不是附加标签，而是本样例中最小的 holonomy 层。
window-\(6\) 的三条 boundary 词
\[
X_6^{\mathrm{bdry}}=\{\texttt{100001},\texttt{100101},\texttt{101001}\}
\]
都具有 dyadic 二层预像对，且纤维内差值固定为 \(34=F_9\)；因此每条 boundary 纤维都天然携带一个规范 \(\ZZ/2\ZZ\) sheet parity。
配合上一节的 boundary 二层覆盖图与 boundary/holonomy 审计表，三条纤维对应的规范 involution 共同生成中心子群
\[
C_{\partial}\cong (\ZZ/2\ZZ)^3,
\]
并在阿贝尔化电荷格
\[
\bigl(G_6^{\mathrm{bin}}\bigr)^{\mathrm{ab}}\cong (\ZZ/2\ZZ)^{21}
\]
中确定一条三维坐标子格。
进一步地，若只允许保持稳定类型标签且由强局域 geometric uplift 实现的翻转，则几何可实现子群仅为对角
\[
P_6^{\mathrm{geo}}=\langle(1,1,1)\rangle\cong \ZZ/2\ZZ,
\]
从而余商
\[
C_{\partial}/P_6^{\mathrm{geo}}\cong (\ZZ/2\ZZ)^2
\]
精确给出两位不可被局域几何吸收的二值自由度。
于是，本样例中的“曲率候选”首先不是连续张量，而是闭路平坦化失败后残留的两位不可约 boundary holonomy。

这一段宜保持“局部—全局障碍”的语义，而不应抢先把上述余商解释成连续场论变量。

\subsubsection{预算裂分：可逆预算与异常预算不是同一个任务}\label{subsubsec:typed-address-biaxial-completion-window6-results-budget-split}
围绕预算裂分的结果段，可写成如下形式。

window-\(6\) 在可逆性与对称保持之间展示出严格的预算裂分。
就裸的精确回放而言，最大纤维顶谱已经决定最小完备补账只需 \(2\) bit；但若进一步要求保持 canonical boundary parity 或 full-center 结构，则局部 completion 的最小开销立即抬升到更高的任务等级。
配合上一节的预算裂分图及预算—脚本映射表，这一分裂可被直接写成
\[
B_6^{\mathrm{replay}}=2,
\qquad
B_6^{\mathrm{bdry}}=3,
\qquad
B_6^{\mathrm{anom}}=21,
\]
因而
\[
B_6^{\mathrm{anom}}-B_6^{\mathrm{replay}}=19.
\]
这说明“可逆恢复”“保持 boundary parity”与“异常签名完备”是三种不同任务等级，不能被压成同一个位数。
在本文的资源语言里，\(L=2\) 与 \(L=21\) 不应被理解成同一预算轴上的大小差异，而应被理解为两种不同语义任务的层级分裂。

这一段最好放在 boundary holonomy 之后，因为只有先看到 \(C_{\partial}\cong (\ZZ/2\ZZ)^3\) 及其对角子群，预算溢价 \(19\) 的语义才不会被误读成技术冗余。

\subsubsection{固定窗统计：表观随机首先表现为闭式可算的非均匀碰撞}\label{subsubsec:typed-address-biaxial-completion-window6-results-fixed-window-statistics}
围绕 fixed-window 统计的结果段，可写成如下形式。

即使只看固定窗统计，window-\(6\) 的可见输出也已经明显偏离均匀基线。
由刚性直方图 \(2{:}8,3{:}4,4{:}9\) 直接计算可得二点碰撞常数
\[
\mathrm{Col}_2^{\mathrm{bin}}(6)=\frac{53}{1024},
\]
而配合上一节的固定窗 collision 图，其相对于均匀基线 \(u_6\) 的二阶偏差满足
\[
\|p_6-u_6\|_2^2=\frac{89}{21504},
\qquad
\chi^2(p_6\|u_6)=\frac{89}{1024}.
\]
因此，本样例中的“像随机”并不是均匀扩散意义下的随机，而是由大纤维抬升的可见统计外观：低阶上表现为非均匀碰撞，高阶上则进一步朝最大纤维冻结。
这也是后续把“确定性粗糙”与“真正噪声”区分开的第一层闭式证书。

这一段的任务，只是先把“非均匀且可审计”锁定下来，而不应过早引入更强的多重分形或物理接口语言。

\subsubsection{粗糙可见量的占位段}\label{subsubsec:typed-address-biaxial-completion-window6-results-rough-visibility-placeholder}
围绕粗糙可见量的结果段，在 `rough-visibility` 审计未完成之前，只能以占位模板出现。
其可用的正文版本如下。

在将可见稳定类型、boundary parity 与动态账本前缀同时注入多尺度相位后，截断可见量
\[
Y_\omega^{(\le J)}(t)
=
\sum_{j=0}^{J} a^j
\cos\!\bigl(
2^jt+\phi_j^{\mathrm{vis}}+\phi_j^{\mathrm{bdry}}+\phi_j^{\mathrm{ledger}}
\bigr)
\]
在 dyadic 放大下表现出持续振荡，其局部振荡代理 \(\operatorname{Osc}_{J,n}\) 与离散 \(p\)-variation 代理 \(V_p^{(J,N)}\) 在待审计参数区间上呈现可复核的尺度分离行为。
配合上一节的粗糙可见量采样图，若与仅保留 \(\phi_j^{\mathrm{vis}}\) 的参考轨道相比，加入 \(\phi_j^{\mathrm{bdry}}\) 与 \(\phi_j^{\mathrm{ledger}}\) 后的轨道在某一稳定尺度区间上出现系统分离，则应把这一分离解释为 residual 相位在多尺度下被重新注入可见层，而不是解释为可见稳定层本身的平凡振荡。

但若数值审计并未显示稳定分离，则该段最后一句必须改写为：
当前读出协议尚未把 residual 有效送回可见层，因此该粗糙可见模型仍需重设相位注入规则。

\subsubsection{机制判别的占位段}\label{subsubsec:typed-address-biaxial-completion-window6-results-mechanism-placeholder}
围绕块碰撞与 mixed-collision 的结果段，同样应在 `block-collision` 审计通过前保持占位状态。
其可用的正文版本如下。

进一步地，块层统计表明机制差异首先出现在尾部几何而不是一体边缘。
对长度为 \(\ell\) 的可见块，块碰撞缺口
\[
\Delta_q^{(\ell)}
:=
\log \mathrm{Col}_q^{(\ell)}-\ell\log \mathrm{Col}_q^{(1)}
\]
若在待审计的 \((q,\ell)\) 区间上显著偏离 \(0\)，则说明该输出并不等价于同边缘的 i.i.d.\ 噪声。
与此同时，归一化 mixed-collision 相似度
\[
\Xi_q^{(0,1)}
:=
\frac{\mathrm{Col}_q^{(0,1)}}
{\bigl(\mathrm{Col}_q^{(0,0)}\mathrm{Col}_q^{(1,1)}\bigr)^{1/2}}
\]
若在 deterministic rough 与某个基线机制的比较中随 \(q\) 增大而出现系统性下降或分叉，则表明两种机制在高阶尾部几何上并不等价。
在这种情形下，本样例中的表观随机应被理解为确定性粗糙系统在有限分辨率下的统计投影，而不是本体层的无结构噪声。

反之，若 \(\Delta_q^{(\ell)}\) 始终接近 \(0\) 且 \(\Xi_q^{(0,1)}\) 没有给出稳定分离，则该段必须被改写为：
当前读出在块层上仍近似可因子化，尚不足以把粗糙机制与噪声基线稳定区分开来。

\subsubsection{跨尺度限制：window-\texorpdfstring{$6$}{6} 的 parity 不是自动外推律}\label{subsubsec:typed-address-biaxial-completion-window6-results-crossscale-limit}
围绕跨尺度限制的结果段，可直接写成如下形式。

window-\(6\) 的 boundary sheet parity 只是一个最小锚点事实，而不是自动跨尺度成立的统一残余律。
具体地，在同一 dyadic uplift 口径下，boundary 差集在 \(m=6\) 时为
\[
\Delta_6^{\mathrm{bdry}}=\{0,34\},
\]
但升到 \(m=7,8\) 时分别变为
\[
\Delta_7^{\mathrm{bdry}}=\{0,55,89\},
\qquad
\Delta_8^{\mathrm{bdry}}=\{0,89,144\},
\]
覆盖层数由二层转为三层。
配合上一节的跨尺度协议门图，这一现象说明：任何关于“sheet parity 跨尺度保真”的叙述，都必须先显式给出并核验一个重锁定协议门
\[
\Gamma_m:\Delta_m^{\mathrm{bdry}}\to \widehat{\Delta}_m^{\mathrm{bdry}}.
\]
在没有 \(\Gamma_m\) 的情况下，window-\(6\) 的 \(\ZZ/2\ZZ\) parity 只能被视为局部锚点，而不能被直接推广为一般分辨率下的无条件残余选择律。

这一段的任务，是主动给出模型边界，阻止最漂亮的 window-\(6\) 事实被误外推成全尺度定理。

\subsubsection{总括段模板}\label{subsubsec:typed-address-biaxial-completion-window6-results-summary-template}
围绕首批结果的总括段，可写成如下形式。

综上，window-\(6\) 已经同时展示出四种彼此独立而又相互咬合的结构：可见层的多壳纤维谱、boundary 扇区的不可平坦化 holonomy、可逆预算与异常预算的严格裂分，以及 fixed-window 层面的非均匀碰撞偏差。
由此，它不再只是一个有限枚举玩具，而应被视为本文的首个最小可审计块：在这个块中，粗糙可见性、表观随机、局部—全局障碍与账本外置第一次被压到同一组有限对象上。
真正的连续极限问题，不在于这些结构是否存在，而在于它们在更大窗口与更细协议下如何被重新锁定、保持或洗出。

\subsubsection{小结}\label{subsubsec:typed-address-biaxial-completion-window6-results-summary}
因此，\S\ref{subsec:typed-address-biaxial-completion-window6-results-paragraph-templates} 的功能，并不是替代上一节的图表说明，而是把图表、表格与协议门进一步压成正文里可复用的结果句法。
在这个句法中，每段只承担一个 claim，evidence 始终回钩到上一节已经锁定的对象图与对象表，而 statistic 则只给出最小且充分的常数证书。
下一步最自然的工作，便是把这些结果模板继续推进为讨论段模板，从而系统区分哪些解释已经获得证据支持，哪些解释仍停留在待验证的外推层。
